\documentclass[12pt,a4paper]{report}

% ---------- Packages ----------
\usepackage[a4paper,margin=1in,headheight=15pt]{geometry}
\usepackage[T1]{fontenc}
\usepackage[utf8]{inputenc}
\usepackage{lmodern}
\usepackage{microtype}
\usepackage{xcolor}
\usepackage{graphicx}
\usepackage{array}
\usepackage{longtable}
\usepackage{booktabs}
\usepackage{enumitem}
\usepackage{fancyhdr}
\usepackage{listings}
\usepackage{hyperref}
\usepackage{titlesec}
\usepackage{float}

\definecolor{codebg}{RGB}{247,248,250}
\definecolor{codeframe}{RGB}{190,195,205}
\definecolor{codeblue}{RGB}{30,70,150}
\definecolor{codegreen}{RGB}{30,115,70}
\definecolor{codered}{RGB}{160,45,45}

\hypersetup{
  colorlinks=true,
  linkcolor=codeblue,
  urlcolor=codeblue,
  pdftitle={Problem Solving and Python Programming Laboratory Manual},
  pdfauthor={Department of Computer Science and Engineering}
}

\lstdefinestyle{pythonstyle}{
  language=Python,
  backgroundcolor=\color{codebg},
  basicstyle=\ttfamily\small,
  keywordstyle=\color{codeblue}\bfseries,
  commentstyle=\color{codegreen},
  stringstyle=\color{codered},
  numbers=left,
  numberstyle=\tiny\color{gray},
  numbersep=8pt,
  frame=single,
  rulecolor=\color{codeframe},
  breaklines=true,
  showstringspaces=false,
  tabsize=4,
  columns=fullflexible,
  keepspaces=true
}
\lstset{style=pythonstyle}

\pagestyle{fancy}
\fancyhf{}
\lhead{\small S1 PSP -- 1131}
\rhead{\small Problem Solving and Python Programming}
\cfoot{\thepage}
\renewcommand{\headrulewidth}{0.4pt}

\titleformat{\chapter}[display]
  {\normalfont\bfseries\centering}
  {\large EXPERIMENT NO. \thechapter}
  {0.5ex}
  {\Large\MakeUppercase}

\setlist[enumerate]{itemsep=2pt,topsep=3pt}
\setlist[itemize]{itemsep=2pt,topsep=3pt}

\newcommand{\experimentdate}{\hfill\textbf{Date:} \rule{3cm}{0.4pt}}
\newcommand{\aim}[1]{\section*{AIM}#1}
\newcommand{\result}[1]{\section*{RESULT}#1}
\newcommand{\courseoutcome}[1]{\section*{RELEVANT COURSE OUTCOME}#1}
\newcommand{\procedure}{\section*{PROCEDURE}}
\newcommand{\questions}[1]{\section*{PRACTICAL RELATED QUESTIONS}#1}
\newcommand{\sampleoutput}{\section*{SAMPLE OUTPUT}}

\begin{document}

% ---------- Cover ----------
\begin{titlepage}
\centering
\vspace*{1.2cm}
{\Large\bfseries STATE BOARD OF TECHNICAL EDUCATION KERALA\par}
\vspace{0.4cm}
{\large Directorate of Technical Education\par}
\vspace{1.4cm}
{\large Diploma in Computer Engineering\par}
{\large Semester I\par}
\vspace{1.3cm}
\rule{\textwidth}{1.2pt}\par
\vspace{0.5cm}
{\Huge\bfseries LABORATORY MANUAL\par}
\vspace{0.5cm}
{\Large\bfseries 1131 -- PROBLEM SOLVING AND\\[0.2cm] PYTHON PROGRAMMING\par}
\vspace{0.5cm}
\rule{\textwidth}{1.2pt}\par
\vspace{1cm}
{\large For Diploma Curriculum Revision 2026\par}
\vfill
{\large Prepared by:\par}
{\large Department of Computer Science and Engineering\par}
\vspace{0.3cm}
{\large Academic Year: \rule{4cm}{0.4pt}\par}
\vspace{1cm}
\end{titlepage}

\pagenumbering{roman}
\chapter*{PREFACE}
This laboratory manual is prepared for first-semester diploma students studying
\textit{Problem Solving and Python Programming (1131)}. It follows the experiment
sequence and course outcomes of the supplied 1131 laboratory syllabus. The format
is based on the reference laboratory manual and includes the aim, objectives,
theory, problem statements, algorithms where appropriate, Python programs,
sample outputs, results, and practice questions.

Students should first understand the problem, identify the input and expected output,
develop the logic, write the program, test it with suitable values, and record the
actual output obtained in the laboratory.

\chapter*{COURSE OUTCOME MATRIX}
\begin{longtable}{>{\centering\arraybackslash}p{1.7cm}p{11cm}}
\toprule
\textbf{CO} & \textbf{Course Outcome}\\
\midrule
CO1 & Acquire foundational logical reasoning skills by designing flowcharts and algorithms and develop Python programs incorporating variables, fundamental data types, and standard input/output operations.\\
CO2 & Optimize program execution through the strategic use of control structures and loop control statements.\\
CO3 & Implement effective data management strategies using strings, lists, tuples, sets, and dictionaries by applying operations and built-in methods to solve programming problems.\\
CO4 & Design and develop modular programs using functions, modules, and packages with proper scope management and argument-passing techniques.\\
\bottomrule
\end{longtable}

\chapter*{GENERAL INSTRUCTIONS}
\begin{enumerate}
\item Read the experiment before attending the laboratory session.
\item Write the experiment number, title, and date clearly in the record.
\item Understand the problem statement before writing the algorithm or program.
\item Type and run the program using the configured Python environment.
\item Check the output with more than one suitable input whenever possible.
\item Correct syntax and logical errors, and record the verified output.
\item Keep the record neat and submit it as instructed by the faculty.
\item Follow laboratory rules and use computers and software responsibly.
\end{enumerate}

\chapter*{INDEX}
\renewcommand{\arraystretch}{1.35}
\begin{longtable}{|>{\centering\arraybackslash}p{0.8cm}|p{7cm}|>{\centering\arraybackslash}p{1.5cm}|>{\centering\arraybackslash}p{1.5cm}|p{1.6cm}|}
\hline
\textbf{Sl. No.} & \textbf{Title of Practical} & \textbf{Mapped CO} & \textbf{Page No.} & \textbf{Date} \\
\hline
1 & Algorithm and Flowchart & CO1 & & \\
\hline
2 & Installation of Python and an IDE & CO1 & & \\
\hline
3 & Introduction to Python Programming & CO1 & & \\
\hline
4 & Decision Control Structures & CO2 & & \\
\hline
5 & Selection-Based Control Structure & CO2 & & \\
\hline
6 & Loop Control Structure & CO2 & & \\
\hline
7 & String Manipulations & CO3 & & \\
\hline
8 & List Operations & CO3 & & \\
\hline
9 & Dictionary Operations & CO3 & & \\
\hline
10 & Set Operations & CO3 & & \\
\hline
11 & Built-in Functions & CO4 & & \\
\hline
12 & User-Defined Functions & CO4 & & \\
\hline
\end{longtable}
\renewcommand{\arraystretch}{1}

\clearpage
\pagenumbering{arabic}

% =========================================================
\chapter{ALGORITHM AND FLOWCHART FOR SIMPLE PROBLEMS}
\experimentdate

\aim{To understand algorithms and flowcharts by designing step-by-step solutions for simple problems before implementing them in Python.}

\section*{OBJECTIVES}
\begin{itemize}
\item Understand the purpose and characteristics of an algorithm.
\item Identify common flowchart symbols.
\item Develop a logical sequence of steps for a problem.
\item Represent the steps using a flowchart.
\end{itemize}

\section*{THEORY}
An \textbf{algorithm} is a finite sequence of clear steps used to solve a problem.
A \textbf{flowchart} is a graphical representation of an algorithm using standard
symbols connected by arrows. An oval usually represents Start/Stop, a parallelogram
represents input/output, a rectangle represents processing, and a diamond represents
a decision.

\courseoutcome{CO1 -- Design algorithms and flowcharts and develop foundational logical reasoning skills.}

\section*{PROBLEM STATEMENTS}
\begin{enumerate}
\item Calculate simple interest from principal, rate, and time.
\item Check whether a given integer is odd or even.
\item Display natural numbers from 1 to a given limit \(N\).
\item Calculate an employee's salary from working days and daily wage.
\end{enumerate}

\section*{ALGORITHM / FLOWCHART}
\subsection*{Problem 1: Simple Interest}
Formula: \(\mathrm{SI}=(P\times R\times T)/100\).
\begin{enumerate}
\item Start.
\item Input principal \(P\), rate \(R\), and time \(T\).
\item Calculate \(\mathrm{SI}=(P\times R\times T)/100\).
\item Display SI.
\item Stop.
\end{enumerate}
\subsection*{Problem 2: Odd or Even}
\begin{enumerate}
\item Start; input integer \(N\).
\item If \(N \bmod 2=0\), display ``Even Number''.
\item Otherwise, display ``Odd Number''.
\item Stop.
\end{enumerate}
\subsection*{Problem 3: Natural Numbers up to \(N\)}
\begin{enumerate}
\item Start; input \(N\); set \(i=1\).
\item While \(i\leq N\), display \(i\) and update \(i=i+1\).
\item Stop.
\end{enumerate}
\subsection*{Problem 4: Employee Salary}
Formula: \(\mathrm{Salary}=\mathrm{WorkingDays}\times\mathrm{DailyWage}\).
\begin{enumerate}
\item Start; input working days and daily wage.
\item Calculate salary.
\item Display salary and stop.
\end{enumerate}

\section*{FLOWCHART RECORD}
Draw a separate flowchart for each problem using standard symbols. Connect the symbols with arrows and label the decision branches Yes/No where required.

\result{Algorithms for simple interest, odd/even checking, printing natural numbers, and employee salary calculation were designed. The corresponding flowcharts should be drawn and verified in the record.}

\questions{
\begin{enumerate}
\item What is an algorithm?
\item What is a flowchart?
\item Which flowchart symbol represents a decision?
\item Why should program logic be planned before coding?
\end{enumerate}}

% =========================================================
\chapter{INSTALLATION OF PYTHON AND AN IDE}
\experimentdate

\aim{To install or verify the Python interpreter and configure an Integrated Development Environment (IDE) for writing and executing Python programs.}

\section*{OBJECTIVES}
\begin{itemize}
\item Understand the Python interpreter and IDE.
\item Verify Python installation.
\item Create and execute a Python file.
\end{itemize}

\section*{THEORY}
The Python interpreter runs Python code. An IDE provides tools such as a code editor,
program execution, syntax highlighting, and debugging. Visual Studio Code (VS Code)
can be used for Python programming after installing the Python extension.

\section*{RESOURCES REQUIRED}
A computer with Windows or Ubuntu, internet access for installation, Python, and an IDE such as VS Code.

\section*{PROCEDURE}
\subsection*{A. Verify or install Python in Windows}
\begin{enumerate}
\item Visit \url{https://www.python.org/downloads/} and download the Windows installer.
\item Run the installer and enable ``Add Python to PATH'' when that option is available.
\item Complete installation.
\item Open Command Prompt and run \texttt{python --version}.
\end{enumerate}
\subsection*{B. Verify or install Python in Ubuntu}
\begin{enumerate}
\item Open Terminal.
\item If Python 3 is not installed, use the package manager as directed by the laboratory instructor.
\item Verify with \texttt{python3 --version}.
\end{enumerate}
\subsection*{C. Configure VS Code}
\begin{enumerate}
\item Download VS Code from \url{https://code.visualstudio.com/}.
\item Install and open VS Code.
\item Open Extensions, search for Python, and install the Python extension.
\item Create a file named \texttt{hello.py} and enter the following code.
\end{enumerate}

\begin{lstlisting}
print("Hello, Python!")
\end{lstlisting}

Run the file using the IDE or the appropriate terminal command for the operating system.

\sampleoutput
\begin{verbatim}
Hello, Python!
\end{verbatim}

\result{Python and the selected IDE were verified or configured, and a simple Python program was executed. Record the version and output obtained on the laboratory computer.}

\questions{
\begin{enumerate}
\item What is a Python interpreter?
\item What is an IDE?
\item What is the purpose of the Python extension in VS Code?
\item What is the extension of a Python source file?
\end{enumerate}}

% =========================================================
\chapter{INTRODUCTION TO PYTHON PROGRAMMING}
\experimentdate

\aim{To write and execute basic Python programs using variables, input/output statements, data types, and arithmetic operators.}

\section*{OBJECTIVES}
\begin{itemize}
\item Use variables to store values.
\item Accept input using \texttt{input()} and display results using \texttt{print()}.
\item Perform arithmetic calculations.
\item Convert input strings to numeric values when necessary.
\end{itemize}

\courseoutcome{CO1 -- Develop Python programs using variables, fundamental data types, and standard input/output operations.}

\section*{THEORY}
Python is a high-level, general-purpose programming language. Variables refer to stored
values. The \texttt{input()} function reads user input as text; \texttt{int()} and
\texttt{float()} convert input to numeric types. The \texttt{print()} function displays
information. Arithmetic operators include \texttt{+}, \texttt{-}, \texttt{*}, \texttt{/},
\texttt{//}, \texttt{\%}, and \texttt{**}.

\section*{PROBLEM 1: TEMPERATURE CONVERSION}
Convert temperature from Celsius to Fahrenheit using
\(F=(C\times 9/5)+32\).

\textbf{Algorithm:} Start; input Celsius \(C\); calculate \(F=(C\times9/5)+32\);
display \(F\); stop.

\begin{lstlisting}
celsius = float(input("Enter temperature in Celsius: "))
fahrenheit = (celsius * 9 / 5) + 32
print("Temperature in Fahrenheit:", fahrenheit)
\end{lstlisting}
\sampleoutput
\begin{verbatim}
Enter temperature in Celsius: 30
Temperature in Fahrenheit: 86.0
\end{verbatim}

\section*{PROBLEM 2: EMPLOYEE SALARY}
Calculate salary using working days multiplied by daily wage.

\begin{lstlisting}
working_days = int(input("Enter number of working days: "))
daily_wage = float(input("Enter daily wage: "))
salary = working_days * daily_wage
print("Monthly Salary:", salary)
\end{lstlisting}
\sampleoutput
\begin{verbatim}
Enter number of working days: 26
Enter daily wage: 850
Monthly Salary: 22100.0
\end{verbatim}

\section*{PROBLEM 3: STUDENT PERCENTAGE}
Calculate total and percentage for five subjects, assuming each subject is out of 100.

\begin{lstlisting}
m1 = float(input("Enter marks of subject 1: "))
m2 = float(input("Enter marks of subject 2: "))
m3 = float(input("Enter marks of subject 3: "))
m4 = float(input("Enter marks of subject 4: "))
m5 = float(input("Enter marks of subject 5: "))
total = m1 + m2 + m3 + m4 + m5
percentage = (total / 500) * 100
print("Total Marks:", total)
print("Percentage:", percentage)
\end{lstlisting}
\sampleoutput
\begin{verbatim}
Enter marks of subject 1: 80
Enter marks of subject 2: 75
Enter marks of subject 3: 90
Enter marks of subject 4: 85
Enter marks of subject 5: 70
Total Marks: 400.0
Percentage: 80.0
\end{verbatim}

\result{The programs for temperature conversion, employee salary, and student percentage were executed and their outputs verified.}

\questions{
\begin{enumerate}
\item What is a variable?
\item What is the difference between \texttt{int()} and \texttt{float()}?
\item What is the purpose of \texttt{input()}?
\item What does the modulus operator return?
\end{enumerate}}

% =========================================================
\chapter{DECISION CONTROL STRUCTURES}
\experimentdate

\aim{To develop Python programs using conditional statements to make decisions based on input values.}

\section*{OBJECTIVES}
Understand \texttt{if}, \texttt{if-else}, and \texttt{if-elif-else}; use relational and logical operators; test programs using different inputs.

\section*{THEORY}
A decision statement selects which block of code runs according to a condition.
Python uses indentation to define the body of a conditional statement.

\courseoutcome{CO2 -- Use control structures to solve programming problems.}

\section*{PROBLEM 1: BUS TICKET CHARGE}
For this exercise, the supplied reference rule is: if the passenger is female or older than 60, set the charge to zero; otherwise retain the entered charge.

\begin{lstlisting}
gender = input("Enter gender (male/female): ")
age = int(input("Enter age: "))
charge = float(input("Enter bus charge: "))

if gender.lower() == "female" or age > 60:
    charge = 0

print("Bus Ticket Charge:", charge)
\end{lstlisting}

\sampleoutput
\begin{verbatim}
Enter gender (male/female): female
Enter age: 35
Enter bus charge: 50
Bus Ticket Charge: 0
\end{verbatim}

\section*{PROBLEM 2: LEAP YEAR CHECKING}
A year is a leap year if it is divisible by 400, or divisible by 4 but not by 100.

\begin{lstlisting}
year = int(input("Enter a year: "))

if year % 400 == 0:
    print("Leap Year")
elif year % 100 == 0:
    print("Not a Leap Year")
elif year % 4 == 0:
    print("Leap Year")
else:
    print("Not a Leap Year")
\end{lstlisting}

\sampleoutput
\begin{verbatim}
Enter a year: 2024
Leap Year
\end{verbatim}

\section*{PROBLEM 3: ELECTRICITY BILL}
Use the rates specified in the reference manual: up to 100 units at Rs. 5 per unit, above 100 and up to 300 at Rs. 8 per unit, and above 300 at Rs. 10 per unit. The program below applies the selected rate to all units consumed.

\begin{lstlisting}
units = int(input("Enter electricity units consumed: "))

if units <= 100:
    bill = units * 5
elif units <= 300:
    bill = units * 8
else:
    bill = units * 10

print("Electricity Bill Amount: Rs.", bill)
\end{lstlisting}

\sampleoutput
\begin{verbatim}
Enter electricity units consumed: 250
Electricity Bill Amount: Rs. 2000
\end{verbatim}

\result{The programs using conditional statements were executed and tested.}

\questions{
\begin{enumerate}
\item What is the purpose of an \texttt{if} statement?
\item What is the difference between \texttt{and} and \texttt{or}?
\item Why is the leap-year condition checked in a particular order?
\item What is the role of indentation in Python?
\end{enumerate}}

% =========================================================
\chapter{SELECTION-BASED CONTROL STRUCTURE}
\experimentdate

\aim{To write Python programs using an if-elif-else ladder and match-case statement for selecting one outcome from multiple alternatives.}

\section*{THEORY}
An \texttt{if-elif-else} ladder checks conditions in order. The \texttt{match-case}
statement selects a case based on the value or pattern being matched and is available
in Python 3.10 and later.

\courseoutcome{CO2 -- Apply control structures to select the appropriate action.}

\section*{PROBLEM 1: STUDENT GRADE}
Use the following grading policy: 90 and above: A; 80--89: B; 70--79: C; 60--69: D; below 60: Fail.

\begin{lstlisting}
mark = int(input("Enter student's mark: "))

if mark >= 90:
    grade = "A"
elif mark >= 80:
    grade = "B"
elif mark >= 70:
    grade = "C"
elif mark >= 60:
    grade = "D"
else:
    grade = "Fail"

print("Grade:", grade)
\end{lstlisting}

\sampleoutput
\begin{verbatim}
Enter student's mark: 85
Grade: B
\end{verbatim}

\section*{PROBLEM 2: DAY NAME USING MATCH-CASE}
\begin{lstlisting}
day = int(input("Enter day number (1-7): "))

match day:
    case 1:
        print("Sunday")
    case 2:
        print("Monday")
    case 3:
        print("Tuesday")
    case 4:
        print("Wednesday")
    case 5:
        print("Thursday")
    case 6:
        print("Friday")
    case 7:
        print("Saturday")
    case _:
        print("Invalid Day Number")
\end{lstlisting}

\sampleoutput
\begin{verbatim}
Enter day number (1-7): 4
Wednesday
\end{verbatim}

\result{The grade evaluation and day-name selection programs were executed and tested with valid and invalid inputs.}

\questions{
\begin{enumerate}
\item When is an if-elif-else ladder useful?
\item What does the underscore case \texttt{case \_} represent?
\item Which Python version introduced \texttt{match-case}?
\end{enumerate}}

% =========================================================
\chapter{LOOP CONTROL STRUCTURE}
\experimentdate

\aim{To write Python programs using loops to repeat operations and solve problems involving sequences and repeated calculations.}

\section*{THEORY}
A loop executes a block of code repeatedly. A \texttt{for} loop is commonly used to
iterate over a sequence or range. A \texttt{while} loop repeats while its condition
remains true. The \texttt{break} statement exits a loop, while \texttt{continue}
skips to the next iteration.

\courseoutcome{CO2 -- Use loops and loop-control statements effectively.}

\section*{PROBLEM 1: PRINT NATURAL NUMBERS UP TO \(N\)}
\begin{lstlisting}
n = int(input("Enter the limit: "))

for i in range(1, n + 1):
    print(i, end=" ")
print()
\end{lstlisting}
\sampleoutput
\begin{verbatim}
Enter the limit: 5
1 2 3 4 5
\end{verbatim}

\section*{PROBLEM 2: SUM OF FIRST \(N\) NATURAL NUMBERS}
\begin{lstlisting}
n = int(input("Enter the limit: "))
total = 0

for i in range(1, n + 1):
    total = total + i

print("Sum:", total)
\end{lstlisting}
\sampleoutput
\begin{verbatim}
Enter the limit: 5
Sum: 15
\end{verbatim}

\section*{PROBLEM 3: COUNT VOWELS IN A WORD}
\begin{lstlisting}
word = input("Enter a word: ")
count = 0

for ch in word.lower():
    if ch in "aeiou":
        count += 1

print("Number of vowels:", count)
\end{lstlisting}
\sampleoutput
\begin{verbatim}
Enter a word: Education
Number of vowels: 5
\end{verbatim}

\result{Loop-based programs were executed and checked using suitable input values.}

\questions{
\begin{enumerate}
\item State one difference between \texttt{for} and \texttt{while}.
\item What is an infinite loop?
\item What does \texttt{break} do?
\item Why must a loop's condition or control variable be updated when required?
\end{enumerate}}

% =========================================================
\chapter{STRING MANIPULATIONS}
\experimentdate

\aim{To perform string operations such as indexing, slicing, concatenation, membership testing, searching, replacing, and formatting.}

\section*{THEORY}
A string is a sequence of characters enclosed in quotes. Strings are immutable:
string methods produce new strings rather than changing the original string in place.
Indexing accesses individual characters, slicing extracts a portion, and methods such
as \texttt{upper()}, \texttt{lower()}, \texttt{replace()}, \texttt{find()},
\texttt{count()}, \texttt{strip()}, and \texttt{split()} help process text.

\courseoutcome{CO3 -- Manage string data using operations and built-in methods.}

\section*{PROBLEM STATEMENT}
Develop a customer feedback program that accepts customer details and performs common string operations.

\begin{lstlisting}
name = input("Enter customer name: ").strip()
customer_id = input("Enter customer ID: ").strip()
email = input("Enter email address: ").strip()
feedback = input("Enter feedback: ").strip()

print("First character:", name[0] if name else "")
print("Last character:", name[-1] if name else "")
print("First five feedback characters:", feedback[:5])
print("Last five feedback characters:", feedback[-5:])

first_name = name.split()[0] if name else ""
username = first_name + customer_id
print("Username:", username)
print("Contains 'good':", "good" in feedback.lower())
print("Uppercase:", feedback.upper())
print("Lowercase:", feedback.lower())
print("Capitalized:", feedback.capitalize())
print("Updated feedback:", feedback.lower().replace("bad", "improve"))
print("Position of 'delivery':", feedback.lower().find("delivery"))
print("Count of 'good':", feedback.lower().count("good"))
print("Words:", feedback.split())
print("Email contains @:", "@" in email)
\end{lstlisting}

\sampleoutput
\begin{verbatim}
Enter customer name: Anu
Enter customer ID: 102
Enter email address: anu@example.com
Enter feedback: Good delivery and good service
First character: A
Last character: u
First five feedback characters: Good 
Last five feedback characters: rvice
Username: Anu102
Contains 'good': True
Uppercase: GOOD DELIVERY AND GOOD SERVICE
Lowercase: good delivery and good service
Capitalized: Good delivery and good service
Updated feedback: good delivery and good service
Position of 'delivery': 5
Count of 'good': 2
Words: ['Good', 'delivery', 'and', 'good', 'service']
Email contains @: True
\end{verbatim}

\result{String indexing, slicing, joining, membership checks, and common string methods were demonstrated.}

\questions{
\begin{enumerate}
\item What is a string?
\item What is meant by string immutability?
\item What is the difference between indexing and slicing?
\item What do \texttt{strip()} and \texttt{split()} do?
\end{enumerate}}

% =========================================================
\chapter{LIST OPERATIONS}
\experimentdate

\aim{To create and manipulate lists using indexing, insertion, deletion, membership testing, slicing, sorting, and list methods.}

\section*{THEORY}
A list stores an ordered collection of items and is written using square brackets.
Lists are mutable, so elements can be added, removed, or changed. Common operations
include \texttt{append()}, \texttt{insert()}, \texttt{remove()}, \texttt{len()},
\texttt{sort()}, \texttt{reverse()}, \texttt{count()}, and \texttt{index()}.

\courseoutcome{CO3 -- Manage sequential data using list operations and built-in methods.}

\section*{PROBLEM STATEMENT}
Create a product inventory and demonstrate list operations.

\begin{lstlisting}
products = ["Rice", "Milk", "Bread", "Tea"]
print("Initial list:", products)
print("First product:", products[0])
print("Last product:", products[-1])

products.append("Sugar")
products.insert(1, "Coffee")
print("After adding products:", products)

products.remove("Bread")
print("After removing Bread:", products)
print("Number of products:", len(products))
print("Is Milk available?", "Milk" in products)
print("Slice [0:3]:", products[0:3])

products.sort()
print("Sorted list:", products)
products.reverse()
print("Reversed list:", products)
print("Count of Rice:", products.count("Rice"))
print("Index of Milk:", products.index("Milk"))
print("Final inventory:", products)
\end{lstlisting}

\sampleoutput
\begin{verbatim}
Initial list: ['Rice', 'Milk', 'Bread', 'Tea']
First product: Rice
Last product: Tea
After adding products: ['Rice', 'Coffee', 'Milk', 'Bread', 'Tea', 'Sugar']
After removing Bread: ['Rice', 'Coffee', 'Milk', 'Tea', 'Sugar']
Number of products: 5
Is Milk available? True
Slice [0:3]: ['Rice', 'Coffee', 'Milk']
Sorted list: ['Coffee', 'Milk', 'Rice', 'Sugar', 'Tea']
Reversed list: ['Tea', 'Sugar', 'Rice', 'Milk', 'Coffee']
Count of Rice: 1
Index of Milk: 3
Final inventory: ['Tea', 'Sugar', 'Rice', 'Milk', 'Coffee']
\end{verbatim}

\result{The list was created and updated using common list operations.}

\questions{
\begin{enumerate}
\item How is a list created?
\item How does \texttt{append()} differ from \texttt{insert()}?
\item What does \texttt{len()} return?
\item Are lists mutable?
\end{enumerate}}

% =========================================================
\chapter{DICTIONARY OPERATIONS}
\experimentdate

\aim{To store and manage data as key-value pairs using Python dictionaries.}

\section*{THEORY}
A dictionary stores information as key-value pairs and is written using curly braces.
Keys must be unique. Dictionaries are mutable, so entries can be added, updated, or
removed. Common methods include \texttt{keys()}, \texttt{values()}, \texttt{items()},
and \texttt{pop()}.

\courseoutcome{CO3 -- Use dictionaries to organize and retrieve structured data.}

\section*{PROBLEM STATEMENT}
Create a product dictionary using product IDs as keys and product names as values.

\begin{lstlisting}
products = {101: "Rice", 102: "Milk", 103: "Tea"}
print("Initial dictionary:", products)
print("Product 102:", products[102])

products[104] = "Sugar"
products[103] = "Green Tea"
print("After add/update:", products)

products.pop(101)
print("After removing product 101:", products)
print("Product IDs:", list(products.keys()))
print("Product names:", list(products.values()))
print("Key-value pairs:", list(products.items()))
print("Is product 102 present?", 102 in products)
print("Number of products:", len(products))
print("Final dictionary:", products)
\end{lstlisting}

\sampleoutput
\begin{verbatim}
Initial dictionary: {101: 'Rice', 102: 'Milk', 103: 'Tea'}
Product 102: Milk
After add/update: {101: 'Rice', 102: 'Milk', 103: 'Green Tea', 104: 'Sugar'}
After removing product 101: {102: 'Milk', 103: 'Green Tea', 104: 'Sugar'}
Product IDs: [102, 103, 104]
Product names: ['Milk', 'Green Tea', 'Sugar']
Key-value pairs: [(102, 'Milk'), (103, 'Green Tea'), (104, 'Sugar')]
Is product 102 present? True
Number of products: 3
Final dictionary: {102: 'Milk', 103: 'Green Tea', 104: 'Sugar'}
\end{verbatim}

\result{Dictionary entries were accessed, added, updated, removed, and displayed using dictionary methods.}

\questions{
\begin{enumerate}
\item What is a key-value pair?
\item Can a dictionary contain duplicate keys?
\item What is the purpose of \texttt{items()}?
\item What happens when an existing key is assigned a new value?
\end{enumerate}}

% =========================================================
\chapter{SET OPERATIONS}
\experimentdate

\aim{To create sets and perform common operations including adding, removing, membership testing, union, intersection, and difference.}

\section*{THEORY}
A set is an unordered collection of distinct elements. Duplicate values are stored only
once. Sets are useful for membership checks and comparing groups of values. The
operators \texttt{|}, \texttt{\&}, and \texttt{-} represent union, intersection, and
difference, respectively.

\courseoutcome{CO3 -- Apply set operations to manage collections of distinct values.}

\section*{PROBLEM STATEMENT}
Maintain library members from two branches and compare their membership lists.

\begin{lstlisting}
branch_a = {"Anu", "Binu", "Ciya", "Dev"}
branch_b = {"Ciya", "Dev", "Esha"}

print("Branch A:", branch_a)
branch_a.add("Fathima")
branch_a.discard("Binu")
print("Updated Branch A:", branch_a)
print("Is Anu registered?", "Anu" in branch_a)
print("Number of members in Branch A:", len(branch_a))

print("Union:", branch_a | branch_b)
print("Intersection:", branch_a & branch_b)
print("Members only in Branch A:", branch_a - branch_b)
\end{lstlisting}

\sampleoutput
\begin{verbatim}
Branch A: {'Anu', 'Binu', 'Ciya', 'Dev'}
Updated Branch A: {'Anu', 'Ciya', 'Dev', 'Fathima'}
Is Anu registered? True
Number of members in Branch A: 4
Union: {'Anu', 'Ciya', 'Dev', 'Esha', 'Fathima'}
Intersection: {'Ciya', 'Dev'}
Members only in Branch A: {'Anu', 'Fathima'}
\end{verbatim}
\textit{Note: The order in which set elements are displayed may differ between runs.}

\result{Set membership and the union, intersection, and difference operations were demonstrated.}

\questions{
\begin{enumerate}
\item How is a set different from a list?
\item Why are duplicate elements not retained in a set?
\item What is the difference between union and intersection?
\item What does the expression \(A-B\) represent?
\end{enumerate}}

% =========================================================
\chapter{BUILT-IN FUNCTIONS}
\experimentdate

\aim{To use suitable Python built-in functions to process values and collections.}

\section*{THEORY}
Built-in functions are available in Python without defining them first. Examples include
\texttt{len()}, \texttt{sum()}, \texttt{min()}, \texttt{max()}, \texttt{sorted()},
\texttt{type()}, \texttt{abs()}, and numeric conversion functions such as \texttt{int()}
and \texttt{float()}. Select a function based on the operation required.

\courseoutcome{CO4 -- Use functions and built-in facilities in Python programs.}

\section*{PROBLEM STATEMENT}
Read marks, then use built-in functions to report basic statistics.

\begin{lstlisting}
marks = [78, 85, 92, 67, 88]

print("Marks:", marks)
print("Number of marks:", len(marks))
print("Total:", sum(marks))
print("Highest mark:", max(marks))
print("Lowest mark:", min(marks))
print("Sorted marks:", sorted(marks))
print("Absolute value of -12:", abs(-12))
print("Data type of marks:", type(marks))
average = sum(marks) / len(marks)
print("Average:", average)
\end{lstlisting}

\sampleoutput
\begin{verbatim}
Marks: [78, 85, 92, 67, 88]
Number of marks: 5
Total: 410
Highest mark: 92
Lowest mark: 67
Sorted marks: [67, 78, 85, 88, 92]
Absolute value of -12: 12
Data type of marks: <class 'list'>
Average: 82.0
\end{verbatim}

\result{Built-in functions were used to find collection length, total, minimum, maximum, sorted values, absolute value, and data type.}

\questions{
\begin{enumerate}
\item What does \texttt{len()} return?
\item What is the difference between \texttt{sort()} and \texttt{sorted()}?
\item What does \texttt{type()} return?
\item How can an average be calculated using \texttt{sum()} and \texttt{len()}?
\end{enumerate}}

% =========================================================
\chapter{USER-DEFINED FUNCTIONS}
\experimentdate

\aim{To define and call functions in Python and use parameters, return values, and local variables to create modular programs.}

\section*{THEORY}
A user-defined function is a named block of reusable code created using the \texttt{def}
keyword. Parameters receive values when the function is called. The \texttt{return}
statement sends a result back to the caller. Functions help divide a large problem into
smaller, understandable parts and reduce repeated code.

\courseoutcome{CO4 -- Develop modular programs using functions, scope management, and argument passing.}

\section*{PROBLEM 1: CALCULATE SIMPLE INTEREST USING A FUNCTION}
\begin{lstlisting}
def simple_interest(principal, rate, time):
    interest = (principal * rate * time) / 100
    return interest

p = float(input("Enter principal: "))
r = float(input("Enter rate of interest: "))
t = float(input("Enter time in years: "))

answer = simple_interest(p, r, t)
print("Simple Interest:", answer)
\end{lstlisting}

\sampleoutput
\begin{verbatim}
Enter principal: 10000
Enter rate of interest: 5
Enter time in years: 2
Simple Interest: 1000.0
\end{verbatim}

\section*{PROBLEM 2: CHECK EVEN OR ODD USING A FUNCTION}
\begin{lstlisting}
def even_or_odd(number):
    if number % 2 == 0:
        return "Even"
    return "Odd"

n = int(input("Enter an integer: "))
print(n, "is", even_or_odd(n))
\end{lstlisting}

\sampleoutput
\begin{verbatim}
Enter an integer: 7
7 is Odd
\end{verbatim}

\section*{PROBLEM 3: FUNCTION WITH DEFAULT ARGUMENT}
\begin{lstlisting}
def greet(name="Student"):
    print("Welcome,", name)

greet()
greet("Anu")
\end{lstlisting}

\sampleoutput
\begin{verbatim}
Welcome, Student
Welcome, Anu
\end{verbatim}

\result{User-defined functions were created and called. Parameters, return values, and a default argument were demonstrated.}

\questions{
\begin{enumerate}
\item Which keyword is used to define a function?
\item What is the purpose of \texttt{return}?
\item What is a parameter?
\item What is a default argument?
\item What is the difference between a local variable and a global variable?
\end{enumerate}}

% ---------- End matter ----------
\chapter*{ADDITIONAL PRACTICE}
\begin{enumerate}
\item Write a program to find the largest of three numbers.
\item Write a program to calculate factorial using a loop.
\item Accept a sentence and count its words.
\item Store five student names in a list and search for a given name.
\item Create a dictionary of student names and marks, then update one student's mark.
\item Create two sets of course participants and display common participants.
\item Write a function that returns the square of a number.
\item Write a function to calculate the average of a list of numbers.
\end{enumerate}

\chapter*{REFERENCES}
\begin{enumerate}
\item \textit{1131 -- Problem Solving and Python Programming Laboratory Manual}, Diploma Curriculum Revision 2026, State Institute of Technical Teachers Training and Research, Kalamassery.
\item Python documentation: \url{https://docs.python.org/3/}
\item Python downloads: \url{https://www.python.org/downloads/}
\item Visual Studio Code: \url{https://code.visualstudio.com/}
\end{enumerate}

\end{document}
